\subsection{Marco Teórico}

\subsubsection{Teoría General de Sistemas}

La Teoría General de Sistemas (TGS), propuesta por Ludwig von Bertalanffy en la década de 1950, plantea que cualquier organización, fenómeno o entidad puede ser comprendida como un sistema compuesto por partes interrelacionadas que funcionan como un todo para lograr una serie de objetivos. Desde esta perspectiva, un sistema no se analiza en forma aislada, sino en términos de sus relaciones internas y su interacción con el entorno \cite{Bertalanffy1968}.

Según la TGS, todo sistema está compuesto por subsistemas interdependientes que interactúan entre sí, y a su vez forma parte de un suprasistema o entorno más amplio. Estos subsistemas intercambian información, recursos o energía a través de entradas (\textit{inputs}) y salidas (\textit{outputs}), y buscan mantener un estado de equilibrio dinámico para adaptarse a cambios en el entorno \cite{Bertalanffy1968}. Los principios de retroalimentación, homeostasis, equifinalidad y sinergia son claves para entender el comportamiento sistémico.

En este contexto, Paralegales como proyecto de Software puede concebirse como un sistema socio-técnico complejo, compuesto por diversos subsistemas funcionales: el \textit{backend}, el \textit{frontend}, la base de datos, la infraestructura en la nube, los usuarios y el equipo de desarrolladores, los stakeholders, etc.

Desde la perspectiva de la TGS, la estrategia de evolución arquitectónica adoptada en este trabajo se fundamenta en la necesidad de reconfigurar las relaciones entre los subsistemas existentes, sin reemplazar el sustrato tecnológico de forma disruptiva, para mejorar la capacidad de adaptación, observación y evolución del conjunto. La introducción de un \textit{Backend-for-Frontend}, la refactorización del \textit{backend} bajo un patrón \textit{Handler--Service--Repository} y la consolidación de un sistema transversal de observabilidad pueden leerse como ajustes del acoplamiento y de los canales de retroalimentación entre subsistemas: reducen el acoplamiento entre el cliente y los servicios sensibles, abren canales de retroalimentación operativa (telemetría) que antes no existían, y aumentan la equifinalidad del sistema al preservar la opción de sustituir componentes individuales sin afectar al resto del conjunto. Este enfoque sistémico permite, además, identificar y atender interdependencias críticas en dos planos complementarios: en el plano organizacional y operativo, la coordinación entre equipos, la automatización del despliegue de software (CI/CD) y la integración de servicios distribuidos; y en el plano del código fuente, las dependencias entre componentes, la calidad de las abstracciones y el grado de acoplamiento entre módulos. Atender ambos planos de forma articulada es lo que permite a Paralegales ganar resiliencia y avanzar hacia un estado de equilibrio dinámico más robusto y sostenible frente a los cambios del entorno organizacional y tecnológico.

\subsubsection{Leyes de la evolución del software}

El estudio sistemático de cómo cambia el software a lo largo del tiempo tiene su origen en el trabajo de \textcite{Lehman1980}, quien, a partir de la observación empírica de sistemas de gran escala en operación, formuló un conjunto de leyes sobre la evolución del software. Su aporte conceptual central consiste en distinguir los programas según su relación con el mundo real: mientras un programa cuya especificación es puramente formal puede considerarse correcto o incorrecto de manera definitiva, un programa que resuelve un problema situado en un entorno real (lo que Lehman denomina un programa de tipo E, por \textit{embedded} en su contexto de uso) nunca alcanza un estado final, porque el propio entorno que pretende servir cambia como consecuencia de su existencia.

De ese planteamiento se derivan dos leyes especialmente relevantes para este trabajo. La primera, la ley del \textit{cambio continuo}, sostiene que un sistema en uso real debe ser adaptado de forma permanente, o bien se vuelve progresivamente menos satisfactorio para sus usuarios: la inacción no preserva el valor del sistema, lo erosiona. La segunda, la ley de la \textit{complejidad creciente}, establece que a medida que un sistema evoluciona su complejidad estructural aumenta, a menos que se realice un trabajo explícito y deliberado para contenerla o reducirla \cite{Lehman1980}. Esta segunda ley es la que confiere estatus teórico a la refactorización: el orden estructural no es un subproducto espontáneo del desarrollo, sino el resultado de una inversión intencional que compite por los mismos recursos que las nuevas funcionalidades.

Leídas en conjunto, ambas leyes describen una tensión difícil de eludir. El sistema debe seguir cambiando para conservar su utilidad, y cada cambio tiende a degradar su estructura, lo que a su vez encarece los cambios siguientes. \textcite{Lehman1980} complementa este cuadro con las leyes de conservación de la familiaridad organizacional y de crecimiento continuo, que vinculan el ritmo sostenible de evolución con la capacidad del equipo para comprender el sistema que modifica. La consecuencia práctica es que la evolución del software se comporta menos como un proyecto con fecha de cierre y más como un proceso continuo de mantenimiento del orden estructural.

Este marco explica teóricamente por qué la intervención descrita en este trabajo adopta la forma que adopta. Paralegales es un sistema en producción, con usuarios activos y desarrollo continuo, es decir, un caso arquetípico de programa de tipo E: no puede congelarse, porque el cambio en su entorno (nuevas exigencias del sector \textit{legaltech}, modificaciones en las APIs de la Rama Judicial, crecimiento de la base de usuarios) seguiría degradando su adecuación aunque el código permaneciera intacto. Tampoco puede rehacerse desde cero, porque una reescritura implicaría suspender la adaptación al entorno durante todo el periodo de transición, precisamente lo que la primera ley advierte que no es sostenible. La refactorización arquitectónica emprendida sobre el \textit{frontend} y la capa de servicios se justifica entonces como el trabajo que la segunda ley señala como imprescindible: una intervención estructural destinada a reducir la complejidad acumulada para que el sistema conserve su capacidad de seguir cambiando. Bajo esta lectura, dicho trabajo no constituye una tarea opcional que compita con el desarrollo de nuevas \textit{features}, sino la condición que hace posible seguir desarrollándolas.

\subsubsection{Teoría de la deuda técnica}

La metáfora de la \textit{deuda técnica} fue introducida por \textcite{Cunningham1992} al describir la experiencia de desarrollo del sistema WyCash. Su formulación original es económica: entregar código que funciona pero que no refleja adecuadamente la comprensión actual del problema equivale a contraer una deuda, y esa deuda devenga intereses en forma de esfuerzo adicional en cada modificación posterior. Mientras la deuda se pague mediante reescrituras oportunas, el préstamo resulta racional; si se deja crecer, los intereses acumulados consumen la capacidad de cambio del equipo hasta volver la deuda impagable. Lo relevante de la metáfora es que reconoce como legítima la decisión de endeudarse, en contraste con una lectura moralizante que interpretaría todo atajo como un simple defecto de calidad.

\textcite{Kruchten2012} llevan la metáfora hacia un marco teórico con capacidad analítica. Su contribución consiste en delimitar el concepto para que resulte operativo: la deuda técnica no es cualquier problema del sistema, sino específicamente aquella propiedad del artefacto de software, invisible para el usuario, que resulta de decisiones de diseño o implementación que compran velocidad en el corto plazo a cambio de un costo diferido. Sobre esa delimitación, los autores organizan lo que denominan el \textit{paisaje} de la deuda técnica y distinguen entre elementos visibles, como los defectos funcionales o las \textit{features} pendientes, y elementos estructuralmente invisibles, como las decisiones de arquitectura, el acoplamiento excesivo o la ausencia de instrumentación. Esta noción de \textit{visibilidad} tiene una consecuencia directa de gestión: la deuda que nadie percibe no se prioriza, de modo que hacerla explícita mediante documentación, indicadores o telemetría es un requisito previo para poder decidir sobre ella.

El marco se completa con una taxonomía ampliamente difundida en la literatura, planteada originalmente por Martin Fowler y recogida en la discusión posterior sobre el tema \cite{Kruchten2012}, que clasifica la deuda según dos ejes: si fue contraída de forma deliberada o inadvertida, y si esa decisión fue prudente o imprudente. La deuda \textit{deliberada y prudente} corresponde a un compromiso consciente, tomado con conocimiento de sus consecuencias y de la intención de pagarlo, típicamente para alcanzar una ventana de mercado. La deuda \textit{inadvertida e imprudente}, en el extremo opuesto, resulta del desconocimiento de las prácticas o del dominio, y suele descubrirse solo cuando los intereses ya son elevados. Entre ambos extremos, la deuda deliberada pero imprudente y la inadvertida pero prudente completan un cuadro que permite razonar sobre la deuda existente en términos de decisión y no de culpa.

Aplicado a Paralegales, este marco describe con precisión la situación de partida. La decisión inicial de apoyarse intensivamente en el SDK de Firebase para resolver autenticación, persistencia y consultas en tiempo real, delegando al cliente parte de la lógica de negocio, respondió a un objetivo explícito de la organización: llegar al mercado con rapidez para validar la propuesta de valor en el sector \textit{legaltech} colombiano. Bajo la taxonomía anterior, se trata de deuda predominantemente deliberada, y prudente en la medida en que compró una ventana de validación que difícilmente se habría obtenido de otro modo. Los intereses de ese préstamo son los cuatro frentes problemáticos diagnosticados previamente: la exposición de lógica sensible en el cliente, la fragmentación entre proveedores, la dependencia parcial de un proveedor único y la ausencia de observabilidad y de aseguramiento de calidad automatizado. Buena parte de esa deuda pertenece a la categoría estructuralmente invisible descrita por \textcite{Kruchten2012}, lo que explica que su costo se manifestara de forma difusa, como lentitud para incorporar cambios, antes que como fallos evidentes.

Desde esta perspectiva teórica, la intervención documentada en este trabajo se interpreta como un \textit{pago selectivo} de deuda y no como una liquidación total. En lugar de reescribir la plataforma, se identifican las partidas de mayor interés, es decir, aquellas cuyo costo diferido crece más rápido y afecta a más cambios futuros, y se amortizan de forma dirigida: la lógica crítica se reubica del lado del servidor, las fronteras entre capas se hacen explícitas y la deuda restante se vuelve visible mediante telemetría estructurada y controles automatizados en el ciclo de vida. La decisión de conservar Firebase Auth y Firestore es coherente con este razonamiento: se trata de deuda cuyos intereses, en la etapa actual del producto, resultan menores que el costo de amortizarla, por lo que la opción racional es mantenerla bajo control y bajo observación en lugar de pagarla de inmediato.

\subsubsection{Modularidad, ocultación de información y valor de las opciones de diseño}

El fundamento teórico de la descomposición modular se remonta a \textcite{Parnas1972}, quien cuestionó el criterio entonces habitual de dividir un sistema siguiendo los pasos de su procesamiento. En su lugar propuso la \textit{ocultación de información} (\textit{information hiding}): cada módulo debe definirse por la decisión de diseño que encapsula y oculta al resto del sistema, de modo que su interfaz revele lo menos posible acerca de su funcionamiento interno. El argumento decisivo de Parnas es de naturaleza económica antes que estética: si los módulos se separan por lo que ocultan, entonces las decisiones susceptibles de cambiar quedan confinadas dentro de una única frontera, y un cambio futuro afecta a un solo módulo en lugar de propagarse por todos los pasos del proceso que lo atraviesan. La modularidad, en esta lectura, es una estrategia de anticipación al cambio.

Poco después, \textcite{Stevens1974} aportaron el vocabulario que convirtió esa intuición en criterio evaluable, al formular el \textit{acoplamiento} y la \textit{cohesión} como propiedades estructurales de un diseño. El acoplamiento describe el grado de dependencia entre módulos distintos, mientras que la cohesión describe el grado en que los elementos internos de un módulo contribuyen a un propósito único. Los autores establecen que un buen diseño tiende a minimizar el primero y maximizar la segunda, y proponen escalas cualitativas que permiten comparar alternativas de descomposición en lugar de discutirlas en abstracto. Esta contribución es la que permite afirmar que una arquitectura es estructuralmente mejor que otra con independencia de que ambas satisfagan los mismos requisitos funcionales.

\textcite{BaldwinClark2000} reformulan la modularidad desde la teoría económica y aportan el concepto que resulta más pertinente para este trabajo. Analizando la industria de la computación, muestran que una arquitectura modular genera \textit{opciones de diseño}: cada módulo cuya interfaz permanece estable constituye una opción, entendida en el sentido financiero del término, es decir, el derecho pero no la obligación de sustituir, mejorar o experimentar con ese módulo en el futuro. El valor de la modularidad no reside únicamente en la claridad del código, sino en la cartera de opciones que crea, y una opción conserva valor incluso si nunca se ejerce, porque lo que se adquiere es la capacidad de decidir más tarde y con mejor información. A esta perspectiva estructural se suma la observación de \textcite{Conway1968}, según la cual cualquier organización que diseña un sistema produce un diseño cuya estructura refleja la estructura de comunicación de la propia organización. La conocida ley de Conway implica que las fronteras técnicas y las fronteras humanas no son independientes: una descomposición modular que no se corresponda con la forma en que el equipo se comunica tenderá a erosionarse.

Estos cuatro aportes fundamentan de manera conjunta las decisiones estructurales adoptadas en Paralegales. El patrón \textit{Handler--Service--Repository} en la capa de servicios es una aplicación directa del criterio de Parnas: el módulo de repositorio existe para ocultar una decisión de diseño concreta, la elección del motor de persistencia, detrás de una interfaz que el resto de la lógica de negocio desconoce. En términos de \textcite{BaldwinClark2000}, esa ocultación crea una opción de diseño sobre la persistencia, y el proyecto ofrece una ilustración inusualmente nítida de la tesis de los autores: la migración de Firestore hacia DynamoDB se inició y posteriormente se descartó, de modo que la opción no se ejerció, y sin embargo la abstracción que la sostiene conservó su valor tanto por la testabilidad que habilita como por mantener abierta la decisión futura. De forma análoga, la organización del \textit{frontend} en \textit{Nuxt Layers} bajo principios de \textit{Domain-Driven Design} busca cohesión por dominio funcional y acoplamiento mínimo entre dominios en el sentido de \textcite{Stevens1974}, mientras que el \textit{Backend-for-Frontend} actúa como frontera explícita que oculta al cliente la heterogeneidad de los proveedores subyacentes y concentra en un único punto las decisiones sobre integración externa.

Finalmente, la ley de Conway ayuda a explicar por qué estas fronteras se definieron como se definieron. En un equipo reducido, donde un mismo desarrollador atraviesa \textit{frontend}, capa de servicios e infraestructura, unas fronteras excesivamente numerosas serían difíciles de sostener; la descomposición adoptada busca por ello coincidir con las unidades de trabajo y los repositorios efectivamente existentes, de modo que la estructura del sistema resulte compatible con la estructura de comunicación de la organización que lo mantiene \cite{Conway1968}.
